\documentclass[10pt,a4paper]{amsart}
\usepackage[margin=19mm]{geometry}
\usepackage{amsmath,amssymb,mathtools}
\usepackage[scr=rsfs,cal=euler]{mathalfa}
\usepackage{xfrac}
\usepackage[unicode,hidelinks,bookmarks=false]{hyperref}
\newcommand{\ie}{i.e.\ }
\newcommand{\cf}{cf.\ }
\newcommand{\resp}{respectively\ }
\newcommand{\Subsection}{\subsection}
\providecommand{\nobreakdash}{\nobreak\hbox{-}}
\setlength{\emergencystretch}{2em}
\multlinegap0pt
\usepackage{bm}
\usepackage{bbm}
\usepackage{dsfont}
\usepackage{xspace}
\usepackage{cancel}
\usepackage{mathtools}
\usepackage{amsthm}
\makeatletter
\@ifundefined{theorem}{\newtheorem{theorem}{Theorem}}{}
\makeatother
\title[Successive vertex orderings of graphs]{Successive vertex orderings of graphs}
\author{Prarthana Agrawal, Abdurrahman Hadi Erturk, Ard A. Louis}
\date{Source: arXiv:2604.07408v2 (CC BY 4.0)}
\begin{document}
\maketitle

\noindent\emph{Source and licence.} Successive vertex orderings of graphs. Version arXiv:2604.07408v2 (CC BY 4.0), distributed under the Creative Commons Attribution 4.0 International licence, as recorded in the arXiv licence metadata for this exact version and shown on the abstract page https://arxiv.org/abs/2604.07408v2. Full rights notice: licenses/arxiv-2604.07408-CC-BY-4.0.txt.\par\smallskip

\noindent\textit{Editorial scope.} Counted unit: the Theorem whose source label is \texttt{thm:derivative} in the article (source0.tex lines 1155--1161, source label \texttt{thm:derivative}), delivered together with the article's own complete original proof of it (source0.tex lines 1163--1338, one \texttt{proof} environment of 176 lines). No mathematics is written, repaired or re-derived: statement and proof are the article's own text line for line. Required context retained uncounted: eq:pr\_bi (lines 290--293). Every definition, display and label of those lines is retained unchanged. Omitted from this package and not redistributed: 1--289, 294--1154, 1162--1162, 1339--2269. Nothing inside a retained excerpt is omitted or altered: every retained body is delivered exactly as the article's lines are written, so no omission receipt of any kind accompanies this package. Citations: the retained excerpts carry no citation command at all, so no bibliography entry of the article is transcribed and no reference is redistributed. Reference adaptation: none was needed and none was made; every reference inside the delivered unit resolves inside it, so no unresolved reference or citation is produced. Printed slips kept literal: no printed word, symbol, hypothesis or number of the counted statement or of its proof was altered. Layout and class: the article's own class is \texttt{article} with packages including authblk, inputenc, fontenc, lmodern, microtype, geometry, amsmath, amssymb, amsthm, mathtools, bm, enumitem, hyperref, cleveref; the wrapper is the lane's pinned \texttt{amsart} editorial wrapper at 10pt on A4, which restates the theorem-like environments the retained text needs and adds the article's own macro definitions that the retained text needs (none), with extra wrapper packages (bm, bbm, dsfont, xspace, cancel, mathtools). The delivered numbering is the unit's own. Assets: no retained span contains a figure environment, a graphics-inclusion command, a PDF-page inclusion, a table or a TikZ picture; no image file is copied into this package, the package declares no asset dependency and no image file is redistributed with it. Licence: version arXiv:2604.07408v2 of this article is distributed under the Creative Commons Attribution 4.0 International licence, as recorded in the arXiv licence metadata for this exact version and shown on the abstract page https://arxiv.org/abs/2604.07408v2. A full rights notice is delivered at licenses/arxiv-2604.07408-CC-BY-4.0.txt. Characters: every retained excerpt is pure ASCII, so the delivered engine, XeLaTeX with its default Latin Modern text font, reports no missing character.
\begin{equation}
\label{eq:pr_bi}
\Pr(B_I)=\frac{a(I)}{n}\,b(I)
\end{equation}

\begin{theorem}\label{thm:derivative}
Let \(A_k\) denote the number of linear orderings \(\pi\) of \(V\) whose set of bad vertices has size exactly \(k\). Define \(F(x):=n!P_G(x)\). Then, for each \(k\ge0\),
\[
A_k=\frac{F^{(k)}(-1)}{k!}
\]
In particular, \(A_0=\sigma(G)\).
\end{theorem}

\begin{proof}
Write
\[
F(x)=n!P_G(x)=\sum_{j\ge0} c_j x^j,
\qquad
c_j:=n!\sum_{\substack{I\subseteq V\\ I\ \mathrm{indep},\,|I|=j}} w(I).
\]

To relate \(F(x)\) to the distribution of bad vertices, we first give a combinatorial interpretation of the coefficients \(c_j\).
For a permutation \(\pi\in\Omega\), let
\[
B(\pi):=\{v\in V:\ \pi\in B_v\}
\]
denote the set of bad vertices in \(\pi\), and write
\(
r(\pi):=|B(\pi)|.
\)
Note that \(B(\pi)\) is always an independent set. Indeed, if
\(u,v\in B(\pi)\) were adjacent, then the definition of \(B_u\)
would imply \(\pi(u)<\pi(v)\), while the definition of \(B_v\)
would imply \(\pi(v)<\pi(u)\), a contradiction.

From \eqref{eq:pr_bi},
\[
\Pr(B_I)=\frac{a(I)}{n}b(I)=w(I).
\]
Recall that \(B_I\) is the event that every vertex of \(I\) is bad. Equivalently,
\[
B_I=\{\pi\in\Omega:\ I\subseteq B(\pi)\}.
\]
Since \(\Omega\) consists of \(n!\) equiprobable orderings, it follows that
\[
n!\Pr(B_I)
=
\#\{\pi\in\Omega:\ I\subseteq B(\pi)\}.
\]
Therefore
\[
n!w(I)
=
\#\{\pi\in\Omega:\ I\subseteq B(\pi)\}.
\]
Since \(B(\pi)\) is independent, every subset of \(B(\pi)\) is independent. Hence the number of independent \(j\)-subsets of \(B(\pi)\) is \(\binom{r(\pi)}{j}\).
Consequently, counting ordered pairs \((\pi,I)\) where
\(\pi\in\Omega\) and \(I\) is an independent \(j\)-subset of
\(B(\pi)\) yields
\begin{equation}\label{eq:cj-sum}
c_j
=
\sum_{\substack{I\subseteq V\\ I\ \mathrm{indep},\,|I|=j}}
\#\{\pi\in\Omega:\ I\subseteq B(\pi)\}
=
\sum_{\pi\in\Omega}\binom{r(\pi)}{j},
\end{equation}
where \(\binom{r(\pi)}{j}=0\) whenever \(r(\pi)<j\).

We now expand \(F(x)\) in powers of \(x+1\). Using the binomial identity
\[
x^j=((x+1)-1)^j
=
\sum_{m=0}^{j}
\binom{j}{m}
(x+1)^m
(-1)^{j-m},
\]
we obtain
\[
F(x)
=
\sum_{j\ge0}
c_j
\sum_{m=0}^{j}
\binom{j}{m}
(-1)^{j-m}
(x+1)^m
=
\sum_{m\ge0}
\left(
\sum_{j\ge m}
c_j
\binom{j}{m}
(-1)^{j-m}
\right)
(x+1)^m.
\]

Since the coefficient of \((x+1)^k\) equals
\(F^{(k)}(-1)/k!\), we have
\begin{equation}\label{eq:deriv-coeff}
\frac{F^{(k)}(-1)}{k!}
=
\sum_{j\ge k}
(-1)^{j-k}
\binom{j}{k}
c_j.
\end{equation}

Substituting \eqref{eq:cj-sum} into
\eqref{eq:deriv-coeff} and exchanging the order of summation gives
\[
\frac{F^{(k)}(-1)}{k!}
=
\sum_{\pi\in\Omega}
\sum_{j\ge k}
(-1)^{j-k}
\binom{j}{k}
\binom{r(\pi)}{j}
=
\sum_{\pi\in\Omega}
H_k\!\bigl(r(\pi)\bigr),
\]
where for integers \(r\ge0\),
\[
H_k(r)
:=
\sum_{j=k}^{r}
(-1)^{j-k}
\binom{j}{k}
\binom{r}{j}.
\]

If \(r<k\), then the sum is empty and \(H_k(r)=0\).
Assume now that \(r\ge k\).
Using the identity
\[
\binom{r}{j}\binom{j}{k}
=
\binom{r}{k}\binom{r-k}{j-k},
\]
we obtain
\[
H_k(r)
=
\binom{r}{k}
\sum_{t=0}^{r-k}
(-1)^t
\binom{r-k}{t}.
\]

If \(r=k\), then the sum consists of the single term \(1\), and hence
\(
H_k(k)=1.
\)
If \(r>k\), then the binomial theorem gives
\[
H_k(r)
=
\binom{r}{k}
\sum_{t=0}^{r-k}
(-1)^t\binom{r-k}{t}
=
\binom{r}{k}(1-1)^{r-k}
=
0.
\]
Thus \(H_k(r)=1\) when \(r=k\), and \(H_k(r)=0\) otherwise.
Hence
\[
H_k\!\bigl(r(\pi)\bigr)
=
\mathbf 1_{\{r(\pi)=k\}}
\]
for every permutation \(\pi\), and hence
\[
\frac{F^{(k)}(-1)}{k!}
=
\sum_{\pi\in\Omega}
\mathbf 1_{\{r(\pi)=k\}}
=
\#\{\pi\in\Omega:\ r(\pi)=k\}
=
A_k.
\]

This completes the proof.
\end{proof}

\end{document}
